\section{Conclusão e Trabalhos Futuros}
\label{sec:conclusion}

Este trabalho propôs e aplicou uma abordagem em dois estágios (segmentação socioeconômica por \textit{K-Means} seguida de comparação intragrupo pelo teste de Mann-Whitney) aos 216 municípios maranhenses com dados no IDEB 2025 (anos iniciais do ensino fundamental, rede municipal). O problema que a motivou é conhecido na literatura: comparar municípios de contextos muito distintos tende a atribuir à escola diferenças que vêm da renda.

Quanto à QP1, a segmentação, feita com a renda \textit{per capita} e a proporção de moradores em pobreza extrema do Censo Demográfico 2022, produziu três grupos (vulnerável, intermediário e de alta renda). Essa partição foi escolhida tanto pelo \textit{Silhouette Score} quanto pela estabilidade \textit{bootstrap}, e as duas variáveis contribuem para ela. Quanto à QP2, nos dois grupos com amostra suficiente para a análise intragrupo, a \textit{formação docente} (proporção de professores com ensino superior e com licenciatura) está associada ao desempenho de forma consistente. As diferenças entre os municípios de alto e baixo IDEB são de 10 a 12 pontos percentuais, com efeito médio. No grupo intermediário, a presença de quadra esportiva também se destaca. O poder preditivo conjunto das variáveis, porém, é baixo.

Para gestores e formuladores de políticas, os resultados apontam a formação docente como fator associado ao desempenho de forma recorrente em municípios de renda baixa e intermediária, o que pode subsidiar a discussão sobre a alocação de recursos educacionais em diferentes perfis socioeconômicos. Para a pesquisa em informática na educação, o estudo mostra o uso do agrupamento como etapa de controle de contexto, validada por critérios de separação, estabilidade e ablação, combinada a testes não paramétricos com tamanho de efeito e correção para comparações múltiplas. O procedimento usa apenas dados públicos e \textit{scripts} abertos e pode ser replicado em outras unidades da federação.

Os trabalhos futuros decorrem das limitações identificadas. O primeiro é incorporar à segmentação dimensões não monetárias, como a escolaridade dos adultos, e comparar o \textit{K-Means} com outros algoritmos de agrupamento. O segundo é replicar o procedimento em outros estados do Nordeste e em outras etapas de ensino, o que permitiria também analisar os municípios de alta renda com amostra suficiente. O terceiro é verificar se a ampliação da formação docente eleva o desempenho, o que exige estudos com desenho longitudinal ou quase experimental, com as próximas edições do IDEB e do Censo Escolar.
